\section{Testpyramiden och de utdelade sviterna}
\label{sec:pyramid}

\subsection{Du skriver inga tester}

Varje testsvit i den här kursen är \textbf{utdelad}. Ditt jobb är att köra dem, läsa dem när de
fallerar, och få dem att passera.

Det är ett medvetet upplägg. Att skriva bra tester är ett eget hantverk, och det kommer i en senare
kurs. Här får du i stället se vad tester är \emph{till för}, på ett projekt där de gör verklig
nytta --- och den som har sett en testsvit hitta en bugg i sin egen kod har en betydligt
konkretare bild av vad ett bra test är än den som ombeds skriva ett först.

\subsection{Varför testning avgör den här kursen}

I tidigare projekt gick det att köra programmet och titta. Här går det inte:

\begin{itemize}
  \item \textbf{Hårdvaran finns inte än.} Den byggs parallellt. Första gången du kan köra mot den
    är \kapref{ch:bringup}, tredje passet från slutet.
  \item \textbf{Hårdvaran är inte din.} När den finns felsöks den samtidigt av någon annan. En
    bugg kan ligga i din kod, i deras VHDL, eller i en kopplingstråd.
  \item \textbf{Det mesta går att köra utan hårdvara.} Det är ingen tröstpremie utan en
    designkonsekvens: sömmarna i \kapref{ch:can} finns just för det.
\end{itemize}

\textbf{Regeln: allt som kan köras på laptopen ska vara grönt innan det körs på hårdvara.}
Bring-upen ska leta efter kopplingsfel, inte efter bitpackningsfel.

\subsection{Pyramidens tre nivåer}

\begin{booktable}{@{}llL@{}}
\thead{Nivå} & \thead{Var} & \thead{Vad} \\ \midrule
Enhetstest & \kapref{ch:frameklassen}, \kapref{ch:interface} & En klass i taget: \code{Frame},
  \code{Stub}. \\
Komponenttest & \kapref{ch:komponent}, \kapref{ch:testning} & Flera delar ihop, utan hårdvara:
  \code{EchoNode} genom stubben; \code{Spi} mot en simulerad registerbank. \\
Integrationstest & \kapref{ch:bringup}, \kapref{ch:analys} & Riktig hårdvara, riktiga ledningar. \\
\end{booktable}

Bredden är en rekommendation om \emph{antal}: många snabba tester längst ned, få dyra längst upp.
Ett projekt med tre snabba tester och en person som kopplar om sladdar är pyramiden upp och ned.

\subsection{Att läsa ett testfall}

Varje utdelat testfall har tre delar:

\begin{cppcode}
TEST(Frame, isValidRejectsTooLargeDlc)
{
    // Arrange.
    driver::can::Frame frame{};
    frame.dlc = 9U;

    // Act & Assert.
    EXPECT_FALSE(driver::can::isValid(frame));
}
\end{cppcode}

Arrange säger vilket utgångsläge kravet gäller i, Act säger vad som ska hända, och
\textbf{Assert är kravet}.

Lägg också märke till värdena: \code{8} och \code{9}, \code{0x7FF} och \code{0x800}. Inte
\code{4} och \code{0x400}. \textbf{Buggar bor på gränserna} --- en jämförelse som ska vara
\code{<=} men skrevs \code{<} fungerar för varenda värde utom ett.

\subsection{Att läsa en fallerande svit}

\begin{console}
FAILED: EXPECT_EQ(frame.data[0], 0xAAU) at test_frame.cpp:42
\end{console}

Arbetsgången, och \textbf{utan att öppna din egen kod i steg 1--3}:

\begin{enumerate}
  \item Läs \textbf{testnamnet}. Det säger vilket krav som är brutet.
  \item Läs \textbf{Arrange}. Den säger under vilka förutsättningar.
  \item Läs \textbf{Assert-raden} som fallerade. Den säger vad som förväntades.
  \item Först därefter: öppna din kod.
\end{enumerate}

Steg 1--3 tar en minut och gör steg 4 kort. Att hoppa direkt till steg 4 är den vanligaste
anledningen till att en enkel bugg tar en timme.

\begin{aside}
\note Den utdelade testsviten är ofta en \textbf{mer exakt kravspecifikation} än prosan i
projektbeskrivningen. Prosan är skriven av en människa och kan vara tvetydig; sviten kompilerar
och kör. Undrar du exakt vad en metod ska returnera i ett gränsfall: öppna testfilen innan du
gissar.
\end{aside}
